\section{Вывод}
В ходе выполнения лабораторной работы было проведено модульное тестирование
разложения функции \(\sin x\) в степенной ряд.
Для этого был выбран достаточный в нашей ситуации набор тестов,
обеспечивающий определенное покрытие возможных случаев.

Также было выполнено тестирование программного модуля для работы с B+ деревьями.
Были определены ключевые точки в алгоритме,
на основе которых сформированы тестовые сценарии.
Последовательность выполнения алгоритма сравнивалась с эталонной,
что позволило выявить возможные расхождения и оценить корректность реализации.

В рамках третьей части задания была сформирована доменная модель на основе заданного текста.
Разработано тестовое покрытие для данной модели, что позволило проверить её соответствие исходным требованиям.

В процессе работы были рассмотрены основные понятия тестирования программного обеспечения,
включая классификацию тестов, модульное тестирование, V-образную модель, методы тестирования,
анализ эквивалентности и регрессионное тестирование.
Изучены принципы работы библиотеки JUnit и её применение для автоматизированного тестирования.

Лабораторная работа позволила закрепить знания в области тестирования программного обеспечения и
получить практический опыт в разработке и применении тестов.
